\subsection{Первый полет: Взлет и посадка}

\textbf{Мотивирующий вопрос:} Какова минимальная последовательность команд для безопасного полета?

Управление дроном похоже на управление автомобилем: сначала нужно завести двигатель (\texttt{arm}), затем начать движение.

\subsubsection*{Основные команды}
\begin{enumerate}
    \item \texttt{arm()} — Запуск моторов (арминг). Винты начинают вращаться на холостых оборотах.
    \item \texttt{takeoff()} — Автоматический взлет на высоту около 1 метра.
    \item \texttt{land()} — Автоматическая посадка.
    \item \texttt{disarm()} — Остановка моторов (дизарминг). Используется после посадки или в экстренных случаях.
\end{enumerate}

\begin{warnbox}[Безопасность]
Всегда используйте блок \texttt{try...finally}, чтобы гарантировать остановку моторов (\texttt{disarm}) даже в случае ошибки в программе или нажатия \texttt{Ctrl+C}.
\end{warnbox}

\begin{codefile}[python]{examples/python/09_drone_takeoff.py}
\end{codefile}

Функция \texttt{time.sleep(N)} приостанавливает выполнение программы на N секунд, давая дрону время на выполнение маневра.

\begin{taskbox}[Прыжок]
Измените программу так, чтобы дрон взлетел, повисел 3 секунды, а затем приземлился.
\end{taskbox}
